\chapter{Modules and \texttt{source}}

\section{What is a module in Hayashi}

In \hayashi{}, a module is simply a \texttt{.hay} file. There is no
package system and no export declarations — any definition in the file
becomes available in the scope of whoever loads it.

This model is intentional: the language is still growing, and the
complexity of a module system with namespaces, visibility rules, and
dependency resolution is added when there is real demand. For now,
\lstinline{source} covers all practical use cases.

\section{The \texttt{source} command}

\lstinline{source} executes a \texttt{.hay} file in the current context.
All definitions from the file become visible from that point on:

\begin{lstlisting}[caption={Loading a module}]
source "lib/statistics.hay"
source "lib/plots.hay"

// functions from statistics.hay and plots.hay are available
let result = compute_gini(df)
\end{lstlisting}

\section{Organising a project}

The typical structure of a \hayashi{} project:

\begin{lstlisting}[style=inline]
project/
  main.hay           -- entry point
  lib/
    data.hay         -- loading and cleaning functions
    estimators.hay   -- custom econometric functions
    tables.hay       -- result formatting
  data/
    raw.csv
    panel.dta
  results/
\end{lstlisting}

\texttt{main.hay}:

\begin{lstlisting}
source "lib/data.hay"
source "lib/estimators.hay"
source "lib/tables.hay"

load "data/panel.dta" as df

let df_clean = prepare(df)
let model    = estimate_fe(df_clean, income ~ educ + exp)

regression_table(model, "results/table1.tex")
\end{lstlisting}

\section{Scope of \texttt{source}}

The file loaded with \lstinline{source} executes in the global scope of
the main file. This means:

\begin{itemize}
  \item Functions defined in the file are globally visible.
  \item Top-level variables in the file are also in global scope.
  \item There are no namespaces: if two files define
        \lstinline{fn compute(...)}, the second overwrites the first.
\end{itemize}

\begin{aviso}
  Since there are no namespaces, avoid generic names in reusable modules.
  Prefer \texttt{gini\_compute} over just \texttt{compute}.
\end{aviso}

\section{Creating a library}

A library file is a \texttt{.hay} file with functions and constants, with
no top-level code that produces side effects on load:

\begin{lstlisting}[caption={lib/utils.hay}]
const PI = 3.14159265358979

fn round_to(x, digits) {
  let factor = 10^digits
  int(x * factor + 0.5) / float(factor)
}

fn between(x, lo, hi) {
  x >= lo and x <= hi
}

fn clamp(x, lo, hi) {
  if x < lo { lo } else if x > hi { hi } else { x }
}
\end{lstlisting}

\begin{lstlisting}[caption={main.hay}]
source "lib/utils.hay"

display round_to(3.14159, 2)    // 3.14
display between(5, 1, 10)       // true
display clamp(15, 0, 10)        // 10
\end{lstlisting}

\section{Roadmap}

In future versions, \hayashi{} will have a module system with:

\begin{itemize}
  \item Explicit namespaces (\texttt{use module::function})
  \item Distribution via a central registry
  \item Dependency declarations in a manifest
\end{itemize}

The transition from \lstinline{source} to that system will be
backwards-compatible: \texttt{.hay} files will continue to work as
modules.
